%------------------------------------------------------------------------------%
% Luis A. D'Afonseca
%
%------------------------------------------------------------------------------%
\documentclass[a4paper,12pt,fleqn]{article}

\usepackage{style_avaliacao}

\ShowAnswers

%------------------------------------------------------------------------------%
\begin{document}

\makeHeader{Integrais e Séries}{Exame -- Turma 4}{19/12/2023}

% 01
%------------------------------------------------------------------------------%
\vspace{\baselineskip}
\exercise{25}
Considerando a função \(f(x) = x\ln(x)\)
\begin{enumerate}[label=\alph*)]
  \item Escreva o polinômio de Taylor de 3\textsuperscript{o} grau centrado em 1
  \item Qual o erro cometido ao aproximar a função pelo polinômio para
  \(-0,5\leq x \leq 1,5\)
\end{enumerate}

\begin{answer}
\end{answer}

% 02
%------------------------------------------------------------------------------%
\vspace{\baselineskip}
\exercise{25}
Calcule a Série de Taylor de \(f(x) = \int\frac{e^x-1}{x}dx\)

\begin{answer}
\end{answer}

% 03 - página 479 Stewart
%------------------------------------------------------------------------------%
\vspace{\baselineskip}
\exercise{25}
Calcule a integral \(\int \frac{\sqrt{9-x^2}}{x^2}dx\)

\begin{answer}
\end{answer}

% 04
%------------------------------------------------------------------------------%
\vspace{\baselineskip}
\exercise{25}
Analise a convergência da série
\(
  \sum_{n=1}^\infty \frac{1}{\left(2n+1\right)^3}
\)

\begin{answer}
\end{answer}

%------------------------------------------------------------------------------%
\end{document}
%------------------------------------------------------------------------------%
